\documentclass[10pt,a4paper]{amsart}
\usepackage[margin=19mm]{geometry}
\usepackage{amsmath,amssymb,mathtools}
\usepackage[scr=rsfs,cal=euler]{mathalfa}
\usepackage{xfrac}
\usepackage[unicode,hidelinks,bookmarks=false]{hyperref}
\newcommand{\ie}{i.e.\ }
\newcommand{\cf}{cf.\ }
\newcommand{\resp}{respectively\ }
\newcommand{\Subsection}{\subsection}
\providecommand{\nobreakdash}{\nobreak\hbox{-}}
\setlength{\emergencystretch}{2em}
\multlinegap0pt
\newcommand{\cobrasPsi}[0]{\mathbf \Psi}
\newcommand{\cobrasPhi}[0]{\mathbf \Phi}
\usepackage{xcolor}
\usepackage{bm}
\newtheorem{proposition}{Proposition}
\title{\huge{\textbf{Uncovering Extreme Event Mechanisms for Prediction and Control with Sensitivity-Balanced Projections }}}
\author{Nicholas Zolman, Sajeda Mokbel, Samuel E. Otto, Steven L. Brunton}
\date{Source: arXiv:2606.05618v1 (CC BY 4.0)}
\begin{document}
\maketitle

\noindent\emph{Source and licence.} \huge{\textbf{Uncovering Extreme Event Mechanisms for Prediction and Control with Sensitivity-Balanced Projections }}. Version arXiv:2606.05618v1 (CC BY 4.0), distributed by arXiv under the Creative Commons Attribution 4.0 International licence, http://creativecommons.org/licenses/by/4.0/, as recorded in the arXiv licence metadata for this exact version and shown on the abstract page https://arxiv.org/abs/2606.05618v1. Full licence notice: \texttt{licenses/arxiv-2606.05618-CC-BY-4.0.txt}.\par\smallskip

\noindent\textit{Editorial scope.} Counted unit: the article's proposition si\_prop:lie\_conv (source0.tex lines 862--881). The article's complete original proof of it is delivered with it (source0.tex lines 883--920, one \texttt{proof} environment of 38 lines). No mathematics is written, repaired or re-derived: the statement and the proof are the article's own lines, in the article's own notation, and no retained line is substituted or re-flowed.  Retained uncounted as the source-local context the proof needs: lines 856--859.  Nothing was omitted from any retained span, so the package carries no omission record.  No retained line contains a citation, so the wrapper carries no bibliography and no bibliography entry.  No retained line contains a figure environment, an image-inclusion command or drawing code, so no image is delivered and the package declares no asset dependency.  Editorial formatting only, in addition: an AMS \texttt{amsart} preamble that restates the article's own declarations and macros used by the retained lines, namely the environments \texttt{proposition} on separate counters, together with the standard packages those lines need.  The retained environments print in this document as Proposition 1 in order of appearance, whereas the article prints the same items with the numbering of its own full document.  The complete native source of the article as distributed in the arXiv e-print for this exact version is bundled unchanged as \texttt{sources/202131/source0.tex}.
\begin{align}
\label{eq:si_local_recon}
    P(\mathbf x) = \mathbf{\tilde x}(\eta) = \sum_k \phi_k(\eta) \langle \psi_k, \mathbf x\rangle    
\end{align}

\begin{proposition} \label{si_prop:lie_conv}
Let a group $G$ act transitively on $\Omega$ and $\mathcal X \subseteq \mathcal{C}(\Omega) \cap \mathcal H$ be invariant to $G$ as above. Suppose that  $\cobrasPhi$ and $\cobrasPsi$ define an operator $P$ as in Eq. ~\ref{eq:si_local_recon} that locally $\epsilon$-reconstructs a neighborhood $U$ of $\eta_{\textit{ref}}$ for every $\mathbf x \in \mathcal X$. If $\mu$ is $G$-invariant, then $\mathbf{x}(\eta)$ is globally $\epsilon$-reconstructed by

\begin{align}
    P_{\textit{global}}(\mathbf{x})(\eta) 
        &= \sum_k \phi_k(\eta_{\textit{ref}})\left(\psi_k \star_G \mathbf x\right)(g_\eta) \\
        &= 
        \cobrasPhi(\eta_{\textit{ref}}) 
        \cdot 
        \left(\cobrasPsi \star_G \mathbf{x} \right)(g_\eta)
\end{align}

\noindent where $\star_G$ is group transformed correlation operator:

$$\left(\psi \star_G \mathbf x\right)(g)
= \int_\Omega \psi(g^{-1} \cdot \eta) \mathbf{x}(\eta)d\mu(\eta)
$$


\end{proposition}

\begin{proof}
Define the function $\mathbf y(\xi) = \mathbf x(g_\eta\cdot  \xi)$ for arbitrary $\xi, \eta \in \Omega$. Since $\mathcal{X}$ is $G$-invariant, $\mathbf y \in \mathcal X$. Choose $\xi = \eta_{\textit{ref}}$, then  $\mathbf y(\eta_{\textit{ref}}) = \mathbf x(g_\eta\cdot  \eta_{\textit{ref}}) = \mathbf x(\eta)$. 
%
Because every signal is locally $\epsilon$-reconstructed in a neighborhood of $\eta_{\textit{ref}}$ and $\eta_{\textit{ref}}$ trivially lives in this neighborhood, we have $|\mathbf y(\eta_{\textit{ref}})  -  P (\mathbf{y})(\eta_{\textit{ref}})| < \epsilon$ where 

$$
 P (\mathbf{y})(\eta_{\textit{ref}})   
    = \sum_k \phi_k(\eta_{\textit{ref}}) \langle \psi_k, \mathbf y\rangle
$$

\noindent Substituting in the definition for the inner product: 
\begin{align*}
    \langle \psi_k, \mathbf y \rangle 
    &= \int_\Omega \psi_k(\xi) \mathbf y(\xi) d\mu(\xi)
    \\
    &= \int_\Omega \psi_k(\xi) \mathbf x(g_\eta \cdot \xi) d\mu(\xi) \\
    &= \int_\Omega \psi_k(g_\eta^{-1} \cdot \zeta) \mathbf x(\zeta) d\mu(\zeta) \\
    &= (\psi_k \star_G \mathbf x)(g_\eta)
\end{align*}

\noindent where we made the substitution $\xi = g_\eta^{-1} \cdot \zeta$ and used the fact that $\mu$ is $G$-invariant: $d\mu(g\cdot \zeta) = d\mu(\zeta)$ for all $g\in G$. Therefore:


$$
P(\mathbf y)(\eta_{\textit{ref}})= \sum_k \phi_k(\eta_{\textit{ref}}) (\psi_k \star_G \mathbf x)(g_\eta) = P_{\textit{global}}(\mathbf{x})(\eta)
$$

Since $\mathbf y(\eta_{\textit{ref}}) = \mathbf{x}(\eta)$, we have 

\begin{align*}
    &|P(\mathbf{y})(\eta_{\textit{ref}}) - \mathbf{y}(\eta_{\textit{ref}})| < \epsilon \\
    \Rightarrow 
    &|P(\mathbf{y})(\eta_{\textit{ref}}) - \mathbf{x}(\eta)| < \epsilon \\
    \Rightarrow
    &|P_{\textit{global}}(\mathbf{x})(\eta) - \mathbf{x}(\eta)| < \epsilon 
\end{align*}

\end{proof}

\end{document}
